\documentclass[11pt]{article}

\usepackage[T1]{fontenc}
\usepackage[margin=1in]{geometry}
\usepackage{booktabs}
\usepackage{tabularx}
\usepackage{listings}
\usepackage{xcolor}

\lstset{
  basicstyle=\ttfamily\small,
  breaklines=true,
  frame=single,
  columns=fullflexible
}

\title{ADR-0002: Production-Grade Test Harness, Automated Reporting, and CI/CD Integration for the LexLean Wireshark Rewrite}
\author{Wireshark LexLean Rewrite Working Group}
\date{2026-11-23}

\begin{document}
\maketitle

\begin{abstract}
This decision record proposes a production-grade verification harness for the LexLean Wireshark
rewrite at \texttt{/run/media/multiplic/AGI/prismagi/mount/SUBSYSTEMS/production}. The harness
follows the operating pattern established in \texttt{kuramoto-prism}: one automated runner, nine
fail-closed suites, generated Markdown and JSON reports, a local pre-commit gate, and a GitHub
Actions workflow that publishes the human-readable report to the job summary. The suite content is
adapted to the current project's Wireshark/pcapng/dissector/filter model and the enterprise capture
and alert model governed by ADR-0001. No Kuramoto numerical invariants or stress claims are copied,
and no formal source is changed by this decision.
\end{abstract}

\section{Context}

ADR-0001 expanded the project to fourteen verified LexLean modules: five Wireshark modules, five
enterprise capture/alert modules, three PRISM facet modules, and the top-level \texttt{Main} module.
The project currently has a single manual verification script, \texttt{verify.sh}, which runs:

\begin{lstlisting}
lexlean fmt --all --check
lexlean lock --check
lexlean check --all
lexlean verify --all
lake build --allow-empty
prismpm check
prismpm build
\end{lstlisting}

That script is sufficient for an ad hoc local check, but it does not produce an audit artifact,
distinguish failing suites, enforce commit-time gates, or publish results to CI. The reference
subsystem \texttt{kuramoto-prism} solves these problems with \texttt{run\_tests.py},
\texttt{run-tests.sh}, \texttt{reports/TEST\_REPORT\_LATEST.md},
\texttt{reports/test\_report\_latest.json}, \texttt{.githooks/pre-commit},
\texttt{install-hooks.sh}, and \texttt{.github/workflows/verify-test.yml}.

Two defects in the reference pattern must not be repeated. First, its report hard-codes Kuramoto
counts and domain-specific invariants; this project must derive counts from its own configuration,
tool output, and source. Second, its CI tool-installation step is a placeholder; this project must
install pinned, digest-checked tools or fail closed.

The implementation plan governing this proposal is
\texttt{docs/PLAN-production-test-harness-ci-cd.md}. The project does not currently contain a
\texttt{.git} directory, so any hook implementation must explicitly initialize a repository on
\texttt{main}.

\section{Decision Drivers}

\begin{enumerate}
  \item \textbf{Single source of truth.} One runner defines suite commands, pass conditions, and
        report content for local, hook, and CI execution.
  \item \textbf{Fail-closed verification.} Any formatter, lockfile, type-check, kernel-verification,
        artifact, or release failure blocks the commit and the CI job.
  \item \textbf{Human-readable audit trail.} Every run produces a Markdown report suitable for
        reviewers and a JSON report suitable for automation.
  \item \textbf{Documentation freshness.} A committed report must correspond to the current formal
        findings; timestamp changes alone must not be treated as semantic drift.
  \item \textbf{Project-specific claims.} Counts, hashes, module names, thresholds, and invariants
        come from this project, not from \texttt{kuramoto-prism}.
  \item \textbf{Reproducible environments.} Tool versions and digests are pinned; the CI workflow
        contains no placeholder installation step.
  \item \textbf{Additive change.} No \texttt{.lex.tex} module, lexicon, formal threshold, or
        ADR-0001 artifact is modified.
\end{enumerate}

\section{Decision}

We establish a nine-suite automated harness, a generated report contract, a local pre-commit gate,
and a GitHub Actions verification and report-publishing pipeline. The runner wraps the existing
verification surface; \texttt{verify.sh} is later reduced to a thin delegate so that command lists
cannot diverge.

\subsection{Suite Catalog}

\begin{table}[htbp]
\centering
\small
\begin{tabularx}{\textwidth}{p{1.25cm} p{3.7cm} X}
\toprule
\textbf{Suite} & \textbf{Name} & \textbf{Command / pass condition} \\
\midrule
\texttt{TS-01}
  & Toolchain Pin \& Runtime Environment
  & \texttt{lean}, \texttt{lake}, \texttt{lexlean}, \texttt{prismpm}; versions must match
    \texttt{lean-toolchain} and the pinned tool manifest. \\
\texttt{TS-02}
  & Canonical Format \& Lock Parity
  & \texttt{lexlean fmt --all --check} and \texttt{lexlean lock --check} must both pass. \\
\texttt{TS-03}
  & Lean Lake Host Build \& Hygiene
  & \texttt{lake build --allow-empty} must exit zero. \\
\texttt{TS-04}
  & Semantic IR AST \& Module Type Check
  & \texttt{lexlean check --all}; parse \texttt{checked 14 modules}, source hash, semantic hash. \\
\texttt{TS-05}
  & Lean 4 Kernel Verification \& Attestation
  & \texttt{lexlean verify --all}; parse \texttt{verified 14 modules} and the 64-hex attestation;
    no sorries or dynamic axioms. \\
\texttt{TS-06}
  & Formal Declaration \& Threshold Audit
  & Source scan over all fourteen modules; declaration families, counts, and ADR-0001 thresholds
    must match the frozen baseline. \\
\texttt{TS-07}
  & Build \& Attestation Artifact Audit
  & Required files under \texttt{.lexlean/}; semantic ID must agree with \texttt{prismpm check}. \\
\texttt{TS-08}
  & PrismPM Model Compilation \& Production Release
  & \texttt{prismpm check}; \texttt{prismpm build}; verify
    \texttt{release/build/HASH/manifest.json} and
    \texttt{model.prism.json}. \\
\texttt{TS-09}
  & Wireshark/Enterprise Invariant Matrix \& Report Integrity
  & Seven local sub-suites described below; all must pass and reports must agree. \\
\bottomrule
\end{tabularx}
\caption{Nine fail-closed verification suites.}
\end{table}

The reference project's \texttt{TS-08} release path is \texttt{production/build/HASH}
because its \texttt{prismpm.toml} uses that build root. This project declares
\texttt{build\_root = "release"}, so the harness must verify
\texttt{release/build/HASH} instead. Copying the reference path literally
would produce a false failure. In these paths, \texttt{HASH} denotes the published 64-hex
PrismPM release identifier.

\subsection{Adapted \texttt{TS-09} Sub-Suites}

\begin{table}[htbp]
\centering
\small
\begin{tabularx}{\textwidth}{p{2.0cm} p{4.2cm} X}
\toprule
\textbf{Sub-suite} & \textbf{Target} & \textbf{Evidence} \\
\midrule
\texttt{SUITE-W1}
  & Wireshark core model
  & Capture limits, pcap/pcapng magic values, 32-bit block alignment, dissector ordering/depth,
    filter evaluation, and cross-module rejection theorems. \\
\texttt{SUITE-E1}
  & Capture and normalization
  & Four-tier ordering, VXLAN/GENEVE termination, worker balancing, buffered transport, schema field
    roles, and $\mathrm{type\_uid} = \mathrm{class\_uid}\times 100 + \mathrm{activity\_id}$. \\
\texttt{SUITE-E2}
  & Threat thresholds
  & Administrative-share/executable rule; $\geq 15$ distinct targets in 10 seconds; $\geq 50$ events
    in 60 seconds; $\geq 500$ events in 300 seconds. \\
\texttt{SUITE-E3}
  & Alert operations
  & Baseline suppression, contextual score threshold, severity mapping, endpoint correlation, and
    SOAR response selection. \\
\texttt{SUITE-X1}
  & Joint integration
  & \texttt{Enterprise.Invariants} joint theorem plus top-level \texttt{Main} integration and
    rejection theorems. \\
\texttt{SUITE-F1}
  & PRISM facets
  & Presence and count of \texttt{Foundation.Arch}, \texttt{Foundation.Qual}, and
    \texttt{Foundation.Sec} declarations. \\
\texttt{SUITE-R1}
  & Report integrity
  & JSON parses; Markdown findings match JSON; latest pair is byte-identical to the timestamped
    pair; committed fingerprint equals fresh fingerprint. \\
\bottomrule
\end{tabularx}
\caption{Project-specific replacement for the Kuramoto stress benchmark.}
\end{table}

These are structural and formal audits. They do not assert numerical convergence, oscillator scaling,
order parameters, or any other Kuramoto-specific property.

\subsection{Report Contract}

The runner emits two paired artifacts on every execution:

\begin{lstlisting}
reports/TEST_REPORT_LATEST.md
reports/test_report_latest.json
reports/TEST_REPORT_<UTC>.md
reports/test_report_<UTC>.json
reports/logs/<suite>_<UTC>.log
\end{lstlisting}

The JSON document uses schema identifier \texttt{lexlean-test-report/1}. It records UTC timestamps,
git commit state, toolchain versions, module/theorem/definition/inductive/structure counts, source
and semantic digests, Lean attestation, per-suite status and duration, the seven \texttt{TS-09}
sub-suite results, PrismPM release hash and path, and the overall pass/fail summary.

A semantic \texttt{fingerprint} is computed as SHA-256 over a canonical projection that omits
volatile fields such as timestamps, durations, and absolute log paths. CI recomputes the fresh
fingerprint and compares it with the committed latest JSON. This detects real report drift without
failing merely because a new run occurred at a different time.

The Markdown report contains an executive summary, suite matrix, all fourteen module rows, the
\texttt{TS-09} sub-suite matrix, attestation and artifact audit, and a failure section with the
failing command and captured stderr. It is appended to \texttt{\$GITHUB\_STEP\_SUMMARY} by the CI
workflow.

Only the two \texttt{*\_LATEST.*} files are committed. Timestamped reports and per-suite logs are
retained as CI artifacts and excluded from git to prevent unbounded repository growth.

\subsection{Local Gate}

The local gate consists of \texttt{.githooks/pre-commit} and \texttt{install-hooks.sh}. The hook runs
the full runner in hook mode, blocks the commit on any failure, and on success stages only:

\begin{lstlisting}
reports/TEST_REPORT_LATEST.md
reports/test_report_latest.json
\end{lstlisting}

The installer initializes \texttt{git init -b main} when no repository exists, sets
\texttt{core.hooksPath} to \texttt{.githooks}, and installs a fallback copy in
\texttt{.git/hooks/pre-commit}. No fast path may skip kernel verification or PrismPM release
checks.

\subsection{CI/CD Pipeline}

The workflow \texttt{.github/workflows/verify-test.yml} runs on pushes and pull requests to
\texttt{main}/\texttt{master} and on manual dispatch. It uses branch/ref concurrency cancellation,
checks out the repository, sets up Python 3.11, installs the pinned Lean toolchain from
\texttt{lean-toolchain}, installs pinned LexLean and PrismPM binaries from a checked-in tool manifest
with SHA-256 verification, and executes \texttt{./run-tests.sh --ci}.

The verification job always uploads \texttt{reports/} as artifact
\texttt{wireshark-lexlean-reports} and always appends the latest Markdown report to the job summary.
A non-zero runner exit fails the job. Branch protection requires this \texttt{verify} check before
merge.

A dependent report-publishing job runs only on successful pushes to \texttt{main}. It republishes the
uploaded report bundle as an immutable, timestamp-named artifact or release asset. It never modifies
formal source, lockfiles, or configuration.

\section{Traceability to Current Artifacts}

\begin{table}[htbp]
\centering
\small
\begin{tabularx}{\textwidth}{p{5.0cm} X}
\toprule
\textbf{Current artifact} & \textbf{Harness coverage} \\
\midrule
\texttt{src/Wireshark/*.lex.tex}
  & \texttt{TS-04}, \texttt{TS-05}, \texttt{SUITE-W1}. \\
\texttt{src/Enterprise/CapturePipeline.lex.tex},
\texttt{src/Enterprise/Normalization.lex.tex}
  & \texttt{SUITE-E1}. \\
\texttt{src/Enterprise/FileServerThreats.lex.tex}
  & \texttt{SUITE-E2}. \\
\texttt{src/Enterprise/AlertOperations.lex.tex}
  & \texttt{SUITE-E3}. \\
\texttt{src/Enterprise/Invariants.lex.tex}, \texttt{src/Main.lex.tex}
  & \texttt{SUITE-X1}. \\
\texttt{src/Foundation/*.lex.tex}
  & \texttt{SUITE-F1}. \\
\texttt{lexlean.toml}, \texttt{lexlean.lock}, \texttt{prismpm.toml}, \texttt{lean-toolchain}
  & \texttt{TS-01}, \texttt{TS-02}, \texttt{TS-07}. \\
\texttt{release/build/HASH/}
  & \texttt{TS-08}. \\
Generated \texttt{reports/*}
  & \texttt{SUITE-R1} and the CI job summary. \\
\bottomrule
\end{tabularx}
\caption{Traceability from the fourteen-module baseline to harness suites.}
\end{table}

The expected baseline is fourteen modules, 52 theorems, 99 definitions, 24 inductive types, and 39
structures, with semantic ID
\texttt{82a79bd033b89dc60f05bd8f58518b2f5d77d6e10c7b839ae02296ab422611b2}, Lean attestation
\texttt{d633598a6e4acdde2a11a1ddc39e35eaeb3d8a2d6d61510c5508b266671b2460}, PrismPM release
\texttt{79af3626858c67b10dad330a8e0a4ece40b728d9cde4072af3c9fe34da4a92b3}, and unchanged certified
Rust export \texttt{a69be6edb75c35ed99b44712e5a8e986b45586515bb78ddaeda13c0b4d9b2a9b}. The harness
must derive these values from tool output and a reviewed baseline manifest, not hard-code them in
report prose. Legitimate model changes require an explicit baseline update in the same change set.

\section{Acceptance Criteria}

\begin{enumerate}
  \item The runner reports nine of nine suites passing and exits zero.
  \item JSON and Markdown reports agree on suite count, statuses, hashes, counts, and fingerprint.
  \item Two consecutive unchanged runs produce identical semantic hash, attestation, release hash,
        counts, and fingerprint; only timestamps and durations differ.
  \item \texttt{verify.sh} delegates to the runner and preserves its successful exit behavior.
  \item The pre-commit hook blocks a failing run and stages only the two latest report files on
        success.
  \item The hook installer succeeds even when no \texttt{.git} directory exists.
  \item CI uses pinned, digest-verified tools with no placeholder installation step.
  \item CI uploads reports on success and failure and appends the Markdown report to the job summary.
  \item The \texttt{verify} status check is required by branch protection.
  \item No \texttt{.lex.tex} source file, closed lexicon, or ADR-0001 modeled threshold is changed.
\end{enumerate}

\section{Consequences}

\subsection{Positive}

\begin{itemize}
  \item Reviewers receive a generated, human-readable verification report with every change.
  \item Failures are attributable to a named suite and include the exact command and captured output.
  \item Local commits and remote CI enforce the same formal and release gates.
  \item Report drift becomes a detectable fingerprint mismatch rather than unnoticed prose.
  \item The harness is project-specific: all counts and invariants come from the fourteen-module
        Wireshark/enterprise model.
\end{itemize}

\subsection{Negative and Mitigations}

\begin{itemize}
  \item Full kernel verification increases commit latency. \textbf{Mitigation:} cache only tool
        installation; do not skip \texttt{TS-05} or \texttt{TS-08}.
  \item Maintaining a baseline manifest adds review overhead. \textbf{Mitigation:} baseline changes
        must be explicit and are displayed as drift in the generated report.
  \item The hook mutates the index by staging reports. \textbf{Mitigation:} stage exactly two named
        latest files and print a clear message; CI remains the authoritative gate.
  \item Tool pinning can fail when upstream releases rotate. \textbf{Mitigation:} store digests in
        \texttt{tools.lock} and update them through a reviewed change that reruns the full harness.
\end{itemize}

\section{Status}

Accepted on 2026-11-23. The nine-suite verification harness, report generation engine, pre-commit hook, and CI pipeline have been implemented and verified. Initial canonical fingerprint recorded:
\begin{itemize}
  \item \textbf{Canonical Fingerprint:} \texttt{71cad005580a4c120936b76efe39954226f2e65a4edeaf81a12f2ec4cead793b}
  \item \textbf{Lean Kernel Attestation:} \texttt{d633598a6e4acdde2a11a1ddc39e35eaeb3d8a2d6d61510c5508b266671b2460}
  \item \textbf{Semantic IR ID:} \texttt{82a79bd033b89dc60f05bd8f58518b2f5d77d6e10c7b839ae02296ab422611b2}
  \item \textbf{PrismPM Release:} \texttt{79af3626858c67b10dad330a8e0a4ece40b728d9cde4072af3c9fe34da4a92b3}
  \item \textbf{Verified Reports:} \texttt{reports/TEST\_REPORT\_LATEST.md} and \texttt{reports/test\_report\_latest.json}
\end{itemize}

\end{document}
